\begin{otherlanguage*}{english}
These pages belong to the ordinary categorical foundations of the theory of derivators. Grothendieck is asking how a large category, and then its categories of diagrams, can be recovered from controlled small data. The recurring difficulty is that a bound on the size of a presentation is not yet an invariant of the object presented. The passage from one to the other requires both a factorization argument and an account of the relations which that argument introduces. Retracts, categories of elements, and changes of indexing category are therefore the central constructions of this part of Chapter~XVIII.

There is a cardinal convention to fix at the outset. A category called \emph{grande devant $\pi$} admits cocones on diagrams with at most $\pi$ arrows. In the customary notation for regular cardinals it is $\pi^+$-filtered. An object called $\pi$-accessible in the manuscript is consequently $\pi^+$-presentable. The parameter $\pi$ itself need not be regular. We retain the manuscript's indices, including its later adjective ``regular'' for a useful cardinal satisfying a particular presentation bound. That adjective is a separate definition, not a statement about the set-theoretic cofinality of $\pi$.

The opening, which continues Proposition~4.5.13 from the preceding pages, compares two bounded-presentation filtrations with the intrinsic filtration $\mathcal M_\pi$. In a cocomplete category the basic calculation is straightforward: a colimit indexed by a category with at most $\pi$ arrows, of $\pi$-accessible objects, is again $\pi$-accessible, because a limit of that size in sets commutes with $\pi^+$-filtered colimits. The converse requires generation. When every object is a sufficiently filtered colimit of objects of a small subcategory, the identity of a presentable object factors through one stage. What follows is a retract of a stage. Lemma~4.5.14 suppresses this last distinction. The commentary gives the corrected implication and an explicit counterexample in sets to its unrestricted wording. Its later application to a subcategory closed under bounded colimits remains valid, since such a subcategory also contains the required idempotent splittings.

The discussion of universal strict epimorphic families has a different purpose. It seeks hypotheses under which one bounded presentation already contains the relations needed for descent. The printed direct-product example in groups is defective even after the accompanying footnote: the two coordinate subgroups of $C_2\times C_2$ do not form a strict epimorphic family. Replacing the product by the free product gives the intended example, including failure of epimorphy after pullback. We distinguish this correction from the unresolved proof transition in Corollary~4.5.16, where a universal epimorphic hypothesis is used as though it were universally strictly epimorphic. The strengthened version has a complete proof. The argument given in these pages does not establish the weaker printed version or the entire deduction of case~C of Proposition~4.5.13.

Section~4.6 builds a small category closed under the colimits needed for reconstruction. The qualitative transfinite construction is sound: at each stage there is a set of diagrams and hence a set of chosen colimits; regularity of the stopping ordinal ensures that every bounded diagram appears early enough to be treated. The numerical estimate for the number of new morphisms is a separate assertion. It is false for an arbitrary small starting category, as the source's own later footnotes suspect. A finitary variety with many binary idempotent operations provides a direct counterexample. Thus the existence of a small reconstruction category must be retained while the claimed sharp bound requires another argument. The normalization of a finite initial cardinal in Corollary~4.6.2 repairs an additional elementary edge case, but does not repair that proof.

The resulting cocomplete recognition theorem is the locally presentable theory: a $\pi_0$-paratopos is, with the stated hypotheses restored, locally $\pi_0^+$-presentable. The presheaf reflection and the small-category presentation explain why suitable filtered colimits commute with bounded limits. This places the discussion alongside Gabriel--Ulmer duality and the systematic accounts of accessible categories, without asserting that Grothendieck was drawing on those accounts here \citep{GabrielUlmer1971,MakkaiPare1989,AdamekRosicky1994}. His immediate reference is SGA~4. One of its assertions, that full faithfulness of a filtered-completion functor forces the original objects to be accessible, fails at the stated generality; the rigid real ring supplies the manuscript's counterexample. The intended Ind application instead follows by evaluation of presheaves.

Exactness does not fix the generating cardinal. To answer footnote~59, we construct, for every prescribed infinite $\pi_0$, a sheaf topos whose only $\pi_0^+$-presentable object is initial, although filtered colimits are exact and a small dense category exists. The construction uses a set of cardinal $\pi_0^{++}$ with the topology of complements of subsets of smaller cardinality. Every nonempty support has an increasing $\pi_0^+$-indexed cover by proper opens; restricting a sheaf to this cover rules out presentability at the proposed rank. This is a supplied development of the question, with a complete proof and no claim of priority. The closing examples of \S4.6 also separate the conditions on images, unions and pushouts of monomorphisms. Pointed sets satisfy all six conditions under discussion and give an example beyond abelian categories and topoi.

Section~4.7 removes the assumption of cocompleteness. Now $\mathcal M=\Ind_{\pi_0}C$, and the primed filtration denotes closure under repeated bounded filtered colimits; its meaning must not be imported from the preceding section. The cardinal equation $\pi^{\pi_0}=\pi$ pays for adjoining all the small witnesses needed to make a subcategory of an indexing category sufficiently filtered. We spell out this closure operation, including the old arrows and cone legs omitted from the printed construction. The resulting comparison of presentation size and presentability is the ordinary version of the modern change-of-cardinal arguments \citep[\S3]{Vasey2020Accessible}.

The retract question which Grothendieck leaves at the centre of this comparison has a positive answer in the literature. Chen's bounded closure theorem, and his formulation of the Makkai--Par\'e retract-removal result, show that iterating the specified bounded filtered colimits does not enlarge the resulting class \citep[Corollaries~4.13 and~4.15]{Chen2021Sifted}. We also give an elementary retract argument which preserves the number of arrows. Consequently $\Filt^\pi=\Filt'^\pi$ at every permitted cardinal. At the initial cardinal the common value is $\Kar(C)$, not necessarily $C$: a two-arrow idempotent category is sufficiently filtered at every rank and realizes a nonrepresentable retract. This endpoint explains several otherwise misleading assertions in the text.

The generalized filtration axioms and their useful cardinals require an additional distinction. A filtered slice is a presentation only when its canonical cocone has the correct colimit. Density supplies that condition, but the axioms printed in 4.7.7 do not force it at every level. We exhibit a filtration of sets whose initial slices are terminal and whose initial realization nevertheless misses every nonempty set. This disproves the bare converse in 4.7.12 and identifies the missing hypothesis. Once usefulness is understood through the realization equivalence, every useful cardinal satisfies the quantitative ``regularity'' requirement of 4.7.13. The commentary proves this with a size-controlled cofinal directed poset, so no example of a useful but nonregular cardinal exists under that interpretation. Conversely, the category of nonempty $\aleph_1$-directed posets has unboundedly many nonuseful manuscript cardinals: every singular $\theta>\aleph_1$ of countable cofinality gives failure of $\theta^+$-accessibility. We provide the complete argument, refining the rank choices in \citet[Example~3.8]{Vasey2020Accessible}. This concerns the accessibility spectrum, rather than the presentability rank of one object \citep{LiebermanRosickyVasey2019Sizes}.

Section~4.8 returns to density and gives the conceptual recognition of accessible categories. The author's initial doubt about enlargement of a dense subcategory is eventually withdrawn, correctly: naturality against the smaller dense category determines the missing components. The characterization of the smallest presentable part as an idempotent completion is likewise valid. Two other assertions need more substantial correction. First, Lemma~4.8.11 omits ``pro'' from the conclusion supplied by SGA~4, I, 8.13.2(b). An accessible finite-limit-preserving functor need only be prorepresentable. Locally constant functions from the Cantor space give an explicit counterexample to representability; preservation of all small limits gives the intended representability theorem. Second, the conjecture in footnote~91 that accessibility of a category and its opposite forces essential smallness is false. Sets and partial injections form a self-dual accessible category with objects of every small cardinality \citep[\S3.2]{Heunen2013L2}; we prove the precise presentability criterion directly.

The small-category theorem requires a more restrained verdict. The argument in these pages does not establish the equality endpoint in 4.8.13. The later manuscript proves a strict cardinal inequality, sufficient for Corollary~4.8.14, and explicitly calls it a weaker result. The continuity result checked in \citet{Tendas2022Continuity} also has a strict threshold. We therefore retain the theorem on essentially small accessible categories while leaving that exact endpoint uncertified; this is not a claim that the endpoint is an established open problem today.

Section~4.9 explains why a diagram of ind-objects need not come from diagrams at individual stages. Coherence can impose a genuine obstruction, even for a one-object category with an involution. We extend the SGA~4 counterexample to every infinite filtering bound using a transfinite signed-chain category. The resulting ind-object has an involution which no stage admits. Under the appropriate colimit hypotheses on $C$, Grothendieck's common-index construction repairs the obstruction. The final cofinality lemmas express that construction through simultaneous factorization and equalization. We restore the missing hypotheses in their printed alternatives and complete the particular proof under way on manuscript page~228.

The later homotopical setting preserves this architecture while replacing sets of maps by mapping spaces and strict diagrams by coherent diagrams. Lurie's recognition theorem characterizes presentable infinity-categories through accessibility, cocompleteness and accessible presheaf localization; his diagram-category theorem supplies the corresponding stability under small diagrams \citep[Theorem~5.5.1.1 and Proposition~5.4.4.3]{Lurie2009HTT}. These are structural counterparts of the constructions here. They explain the relevance of this long ordinary-categorical preparation to the work on derivators without identifying strict diagrams in a homotopy category with coherent diagrams in a homotopy theory.
\end{otherlanguage*}
